%============================================================
% BAB I — PENDAHULUAN
% Struktur: A. Latar Belakang Masalah, B. Rumusan Masalah,
% C. Batasan Masalah, D. Tujuan Penelitian, E. Manfaat Penelitian.
% Spasi isi: 2 (termasuk judul bab). Sub-bab pakai huruf A., B., ...
%============================================================
\chapter{PENDAHULUAN}

\section{Latar Belakang Masalah}
% Ganti teks panduan di bawah ini dengan latar belakang Anda.
Bagian ini memuat penjelasan tentang fenomena umum yang terjadi
dan kemudian dipersempit mengarah ke permasalahan yang akan
diteliti atau dibahas. Latar belakang masalah merupakan uraian
yang komprehensif mengenai pentingnya permasalahan tersebut
diangkat menjadi topik penulisan, ditinjau dari aspek praktis
maupun kontribusi ilmiah secara teoritis
\parencite{winowatan2009}.

Penulisan latar belakang masalah disajikan dalam bentuk uraian
yang secara kronologis diarahkan untuk langsung menuju rumusan
masalah. Istilah asing ditulis miring, misalnya \emph{machine
learning} dan \emph{internet of things}. Semua simbol matematika,
variabel, dan persamaan wajib menggunakan sintaks mode matematika
LaTeX, misalnya $x$ dan $y$, atau tampilan rumus seperti pada
persamaan~\eqref{eq:contoh-bab2}.

\section{Rumusan Masalah}
Bagian ini memuat penjelasan tentang permasalahan sehingga
memerlukan solusi. Permasalahan yang diuraikan dalam latar
belakang masalah dirumuskan kembali secara tegas dan jelas dalam
bentuk poin-poin yang terinci yang berisi masalah-masalah yang
akan dikaji. Dirumuskan dalam bentuk pertanyaan atau pernyataan;
jumlahnya dapat satu atau lebih, tergantung kedalaman dan luasnya
masalah yang akan diteliti. Sebagai contoh:
\begin{enumerate}
  \item Bagaimana merancang dan membangun sistem yang diusulkan?
  \item Bagaimana menguji kelayakan sistem yang diusulkan?
\end{enumerate}

\section{Batasan Masalah}
Dalam penulisan batasan masalah harus dihindari penggunaan kata
``tidak dibahas/membahas''. Yang harus dituliskan adalah segala
sesuatu yang menjadi topik pembahasan skripsi. Misalnya:
``Penelitian ini dibatasi oleh hal-hal sebagai berikut:''
\begin{enumerate}
  \item Batasan-batasan permasalahan yang akan dicari solusinya;
  \item Asumsi-asumsi yang digunakan sebagai dasar penulisan;
  \item Batasan solusi.
\end{enumerate}

\section{Tujuan Penelitian}
Penelitian dapat bertujuan untuk mengetahui, menguraikan,
menerangkan, membuktikan, atau menerapkan model, konsep, atau
dugaan. Tujuan adalah kurang lebih dari jawaban/capaian yang ada
di rumusan masalah. Bagian ini memuat penjelasan secara spesifik:
\begin{enumerate}
  \item Solusi yang ditawarkan untuk menangani permasalahan yang dihadapi;
  \item Hal-hal yang ingin dicapai melalui kegiatan yang akan dilakukan.
\end{enumerate}

\section{Manfaat Penelitian}
Manfaat penelitian menguraikan dan menjelaskan kegunaan secara
teoritik dan aplikatif dari penelitian yang telah dilakukan.
Manfaat atau kontribusi penelitian memuat dua hal yang mendasar:
\begin{itemize}
  \item \textbf{Manfaat Teoritik}, apabila hasil penelitian akan
        menghasilkan sebuah pendapat baru atau hasil penerapan hukum;
  \item \textbf{Manfaat Aplikatif}, apabila terdapat manfaat atau
        nilai guna hasil penelitian bagi stakeholder atau pihak-pihak
        yang terkait langsung dengan hasil penelitian.
\end{itemize}
